\chapter{Computational tables and reference data}
\label{app:computational-tables}

\section{Configuration space weight tables}

\subsection{Low-degree Kontsevich weights}

The following table lists the Kontsevich graph weights $w_\Gamma$ appearing in the formality morphism $U \colon T_{\mathrm{poly}} \to D_{\mathrm{poly}}$ on $\bR^d$~\cite{Kontsevich03}.  Each graph $\Gamma$ has aerial vertices (in the upper half-plane) and boundary vertices (on $\bR$); the weight is the integral over the configuration of aerial vertices of a product of angle forms.

\begin{center}
\begin{tabular}{|c|c|c|c|}
\hline
Graph Type & Vertices & Edges & Weight $w_\Gamma$ \\
\hline
Single point & 1 & 0 & 1 \\
\hline
Binary tree & 2 & 1 & 1 \\
\hline
Wheel & 2 & 2 & $\frac{1}{12}$ \\
\hline
Chain & 3 & 2 & $\frac{1}{24}$ \\
\hline
Complete & 3 & 3 & $\frac{\zeta(3)}{(2\pi)^2}$ \\
\hline
\end{tabular}
\end{center}

\section{Affine Kac--Moody data}

\subsection{Classical simple Lie algebras}

\begin{center}
\begin{tabular}{|c|c|c|c|c|}
\hline
Type & $\mathfrak{g}$ & $\dim \mathfrak{g}$ & $h^\vee$ & Level shift $k'$ \\
\hline
$A_n$ & $\mathfrak{sl}_{n+1}$ & $n(n+2)$ & $n+1$ & $-k - 2(n+1)$ \\
\hline
$B_n$ & $\mathfrak{so}_{2n+1}$ & $n(2n+1)$ & $2n-1$ & $-k - 2(2n-1)$ \\
\hline
$C_n$ & $\mathfrak{sp}_{2n}$ & $n(2n+1)$ & $n+1$ & $-k - 2(n+1)$ \\
\hline
$D_n$ & $\mathfrak{so}_{2n}$ & $n(2n-1)$ & $2n-2$ & $-k - 2(2n-2)$ \\
\hline
\end{tabular}
\end{center}

\subsection{Exceptional Lie algebras}

\begin{center}
\begin{tabular}{|c|c|c|c|}
\hline
Type & $\dim$ & $h^\vee$ & Level shift \\
\hline
$G_2$ & 14 & 4 & $-k-8$ \\
\hline
$F_4$ & 52 & 9 & $-k-18$ \\
\hline
$E_6$ & 78 & 12 & $-k-24$ \\
\hline
$E_7$ & 133 & 18 & $-k-36$ \\
\hline
$E_8$ & 248 & 30 & $-k-60$ \\
\hline
\end{tabular}
\end{center}

\section{W-algebra structure constants}

\subsection{\texorpdfstring{$W_3$ commutators (explicit)}{W 3 commutators (explicit)}}

For the $W_3$ algebra with generators $\{L_n\}_{n \in \mathbb{Z}}$ (Virasoro) and $\{W_n\}_{n \in \mathbb{Z}}$ (weight-3 field):

\begin{align}
[L_m, L_n] &= (m-n)L_{m+n} + \frac{c}{12}m(m^2-1)\delta_{m+n,0} \\
[L_m, W_n] &= (2m-n)W_{m+n} \\
[W_m, W_n] &= \frac{c}{360}m(m^2-1)(m^2-4)\delta_{m+n,0} \\
&\quad + \frac{16(m-n)}{22+5c}\Lambda_{m+n} \\
&\quad + (m-n)\frac{2m^2 - mn + 2n^2 - 8}{30}L_{m+n}
\end{align}

\subsection{OPE residue formulas}

For primary fields $\phi_i$ of weight $h_i$, the residue at collision divisor $D_{ij}$ is:
$$\text{Res}_{D_{ij}}[\phi_i(z) \otimes \phi_j(w) \otimes \eta_{ij}] = C^k_{ij} \phi_k(w)$$
where $C^k_{ij}$ is the OPE coefficient at the simple pole.  The residue extracts the $(z-w)^{-1}$ term of the OPE, which contributes when:
$$h_i + h_j - h_k = 1$$
(simple pole condition).  Higher-order poles ($h_i + h_j - h_k > 1$) contribute to higher-order terms in the bar differential.

\section{Arnold relation expansions}

\subsection{Three-point relations}

$$\eta_{12} \wedge \eta_{23} + \eta_{23} \wedge \eta_{31} + \eta_{31} \wedge \eta_{12} = 0$$

\subsection{Four-point relations}

For $n = 4$, there are four Arnold relations, one for each $3$-element subset of
$\{1,2,3,4\}$:
\begin{align}
\eta_{12} \wedge \eta_{23} + \eta_{23} \wedge \eta_{31} + \eta_{31} \wedge \eta_{12} &= 0 \\
\eta_{12} \wedge \eta_{24} + \eta_{24} \wedge \eta_{41} + \eta_{41} \wedge \eta_{12} &= 0 \\
\eta_{13} \wedge \eta_{34} + \eta_{34} \wedge \eta_{41} + \eta_{41} \wedge \eta_{13} &= 0 \\
\eta_{23} \wedge \eta_{34} + \eta_{34} \wedge \eta_{42} + \eta_{42} \wedge \eta_{23} &= 0
\end{align}
These are all instances of the three-point pattern applied to subsets.
Products of \emph{disjoint} pairs $\eta_{ij} \wedge \eta_{kl}$ with
$\{i,j\} \cap \{k,l\} = \emptyset$ are \emph{not} constrained by Arnold relations
and represent independent classes in $H^2(C_4(\mathbb{C}))$

\subsection{\texorpdfstring{General $n$-point relations}{General n-point relations}}

For any three distinct indices $i, j, k$, the fundamental Arnold relation is:
$$\eta_{ij} \wedge \eta_{jk} + \eta_{jk} \wedge \eta_{ki} + \eta_{ki} \wedge \eta_{ij} = 0$$
More generally, the Orlik--Solomon relations state: for each circuit (minimal dependent set) $\{h_0, \ldots, h_p\}$ in the arrangement, one has $\sum_{r=0}^{p} (-1)^r \, e_{h_0} \wedge \cdots \wedge \widehat{e_{h_r}} \wedge \cdots \wedge e_{h_p} = 0$.

\section{Modular forms at higher genus}

\subsection{Genus 1: Eisenstein series}

$$E_2(\tau) = 1 - 24\sum_{n=1}^\infty \frac{nq^n}{1-q^n}, \quad q = e^{2\pi i \tau}$$
$$E_4(\tau) = 1 + 240\sum_{n=1}^\infty \frac{n^3q^n}{1-q^n}$$
$$E_6(\tau) = 1 - 504\sum_{n=1}^\infty \frac{n^5q^n}{1-q^n}$$

\subsection{Genus 2: Siegel modular forms}

For period matrix $\Omega \in \mathbb{H}_2$ and characteristic $\epsilon = (\epsilon', \epsilon'') \in \{0, \tfrac{1}{2}\}^4$ with $\epsilon', \epsilon'' \in \{0, \tfrac{1}{2}\}^2$:
$$\Theta[\epsilon](\Omega) = \sum_{n \in \mathbb{Z}^2} \exp\bigl[\pi i(n+\epsilon')^T \Omega (n+\epsilon') + 2\pi i(n+\epsilon')^T \epsilon''\bigr]$$
(cf.\ the genus-$g$ formula in Appendix~\ref{app:theta}).  The $2^4 = 16$ characteristics split into $6$ odd and $10$ even.

\section{\texorpdfstring{$W_3$ algebra coefficients}{W-3 algebra coefficients}}
\label{app:w3-coefficients}
\label{app:w3-jacobi-full}

\subsection{\texorpdfstring{Composite field $\Lambda$ formula}{Composite field formula}}

\begin{equation}
\Lambda(w) = {:}T(w)T(w){:}
\end{equation}
the normal-ordered product.  The quasi-primary composite appearing in the OPE (cf.\ Example~\ref{ex:w3-completion-full}) is $\Lambda_{\mathrm{qp}} = {:}TT{:} - \tfrac{3}{10}\partial^2 T$; in the mode commutator below we use the \emph{non-quasi-primary} convention $\Lambda_n = \sum_m {:}L_m L_{n-m}{:}$ (without derivative correction), which differs from the quasi-primary normalization used in~\cite{BoSch93} by the $-\tfrac{3}{10}\partial^2 T$ term.

\subsection{Coefficient values for standard central charges}

\begin{center}
\begin{tabular}{|c|c|c|c|c|}
\hline
$c$ & $\alpha(c) = \frac{16}{22+5c}$ & $\beta(c) = \frac{3}{10}$ & 
$\alpha$ (fraction) & $\alpha$ (decimal) \\
\hline
$-2$ & $\frac{16}{12}$ & $\frac{3}{10}$ & $\frac{4}{3}$ & $1.333$ \\
\hline
$-1$ & $\frac{16}{17}$ & $\frac{3}{10}$ & $\frac{16}{17}$ & $0.941$ \\
\hline
$0$ & $\frac{16}{22}$ & $\frac{3}{10}$ & $\frac{8}{11}$ & $0.727$ \\
\hline
$1$ & $\frac{16}{27}$ & $\frac{3}{10}$ & $\frac{16}{27}$ & $0.593$ \\
\hline
$2$ & $\frac{16}{32}$ & $\frac{3}{10}$ & $\frac{1}{2}$ & $0.500$ \\
\hline
$5$ & $\frac{16}{47}$ & $\frac{3}{10}$ & $\frac{16}{47}$ & $0.340$ \\
\hline
$10$ & $\frac{16}{72}$ & $\frac{3}{10}$ & $\frac{2}{9}$ & $0.222$ \\
\hline
$50$ & $\frac{16}{272}$ & $\frac{3}{10}$ & $\frac{1}{17}$ & $0.059$ \\
\hline
$100$ & $\frac{16}{522}$ & $\frac{3}{10}$ & $\frac{8}{261}$ & $0.031$ \\
\hline
\end{tabular}
\end{center}

\subsection{Mode expansion formulas}

The normal-ordered composite $\Lambda$ has mode expansion:
$$\Lambda_n = \sum_{m \in \mathbb{Z}} :L_m L_{n-m}:$$

\emph{Low mode examples.}
\begin{align}
\Lambda_0 &= L_0^2 + 2\sum_{k>0} L_{-k}L_k \\
\Lambda_1 &= 2L_0 L_1 + 2\sum_{k>0} L_{-k}L_{k+1} \\
\Lambda_{-1} &= 2L_0 L_{-1} + 2\sum_{k>0} L_k L_{-k-1}
\end{align}

\subsection{Comparison with literature}

\emph{Convention note.}  Much of the literature (Zamolodchikov 1985, Bouwknegt--Schoutens 1993, Arakawa 2017) defines a ``full composite'' $\widetilde{\Lambda} = \frac{16}{22+5c}\bigl({:}TT{:} - \frac{3}{10}\partial^2 T\bigr)$, absorbing the coefficient $\alpha = 16/(22+5c)$ into the quasi-primary field.  With that convention the commutator reads $(m-n)\widetilde{\Lambda}_{m+n}$ (no extra factor of $\alpha$).  We use $\Lambda = {:}TT{:}$ as above, so our commutator carries the explicit factor $\frac{16(m-n)}{22+5c}\Lambda_{m+n}$.  Both conventions yield the same mode algebra; the quasi-primary version $\Lambda_{\mathrm{qp}} = {:}TT{:} - \frac{3}{10}\partial^2 T$ (cf.\ Example~\ref{ex:w3-completion-full}) is the field-theoretic standard for the OPE.

All sources agree after accounting for normalization differences.
